\documentclass[../../main.tex]{subfiles}% 注意这里的文件路径不能用 ./main.tex ,否则用latexmk编译子文件会报错
\graphicspath{{\subfix{./image/}}} % 指定图片目录,后续可以直接使用图片文件名
% 注意这里的文件路径不能用 ../../image/ ,否则用latexmk编译子文件会报错

% 例如:
% \begin{figure}[H]
% \centering
% \includegraphics[scale=0.3]{图.png}
% \caption{}
% \label{figure:图}
% \end{figure}
% 注意:上述\label{}一定要放在\caption{}之后,否则引用图片序号会只会显示??.

\begin{document}

\section{可分多项式$\,\,$完备域}

\begin{definition}
设$F$为域，$f(x)\in F[x]$，且
\begin{align*}
f(x)=a_nx^n+a_{n-1}x^{n-1}+\dots+a_0,
\end{align*}
定义$F[x]$中多项式
\begin{align*}
f'(x)=na_nx^{n-1}+(n-1)a_{n-1}x^{n-2}+\dots+a_1
\end{align*}
为$f(x)$的\textbf{形式微商}（简称微商）。
\end{definition}

\begin{remark}
形式微商满足如下性质：
\begin{enumerate}[(1)]
\item $a'=0,\forall a\in F$；当$\mathrm{ch}F=0$时，$f'(x)=0$当且仅当$f(x)\in F$.
\item $x'=1$.
\item $(f(x)+g(x))'=f'(x)+g'(x)$.
\item $(cf(x))'=cf'(x),\forall c\in F$.
\item $(f(x)g(x))'=f'(x)g(x)+f(x)g'(x)$.
\end{enumerate}
\end{remark}

\begin{theorem}
设$K$是$f(x)\in F[x]$的分裂域，$\alpha\in K$是$f(x)$的$k$重根，则
\begin{enumerate}[(1)]
\item 若$\mathrm{ch}F\nmid k$，则$\alpha$是$f'(x)$的$k-1$重根；
\item 若$\mathrm{ch}F\mid k$，则$\alpha$是$f'(x)$的至少$k$重根。
\end{enumerate}
\end{theorem}

\begin{proof}
在$K[x]$中，设$f(x)=(x-\alpha)^kg(x)$，其中$g(\alpha)\neq0$。由微商的乘积法则：
\begin{align*}
f'(x)&=k(x-\alpha)^{k-1}g(x)+(x-\alpha)^kg'(x) \\
&=(x-\alpha)^{k-1}\big(kg(x)+(x-\alpha)g'(x)\big).
\end{align*}
当$\mathrm{ch}F\nmid k$时，$k\neq0$，且$kg(\alpha)+(\alpha-\alpha)g'(\alpha)=kg(\alpha)\neq0$，故$\alpha$是$f'(x)$的$k-1$重根。

当$\mathrm{ch}F\mid k$时，$k=0$，于是$f'(x)=(x-\alpha)^kg'(x)$，故$\alpha$是$f'(x)$的至少$k$重根.

\end{proof}

\begin{theorem}
设$K$是$f(x)\in F[x]$的分裂域，则$f(x)$在$K$中无重根的充分必要条件是$(f(x),f'(x))=1$。
\end{theorem}

\begin{proof}
必要性：若$f(x)$在$K$中无重根，则
\begin{align*}
f(x)=(x-\alpha_1)(x-\alpha_2)\dots(x-\alpha_n),
\end{align*}
其中$n=\deg f(x)$，且$i\neq j$时$\alpha_i\neq\alpha_j$。由上一定理，每个$\alpha_i$都不是$f'(x)$的根，故$(f(x),f'(x))=1$。

充分性：若$f(x)$在$K$中有重根$\alpha$，重数$k>1$，则$\alpha$是$f'(x)$的至少$k-1$重根，于是$(x-\alpha)^{k-1}\mid f(x)$且$(x-\alpha)^{k-1}\mid f'(x)$，即$\deg(f(x),f'(x))\geqslant1$，故$(f(x),f'(x))\neq1$.

\end{proof}

\begin{corollary}
设$p(x)$是$F[x]$中的不可约多项式，则$p(x)$在其分裂域中有重根的充分必要条件是$p'(x)=0$。
\end{corollary}

\begin{proof}
由上一定理，$p(x)$有重根当且仅当$(p(x),p'(x))=d(x)\neq1$。
因$p(x)$不可约，$d(x)\mid p(x)$且$\deg d(x)\geqslant1$，故$\deg d(x)=\deg p(x)$。
又$\deg p'(x)<\deg p(x)$，$d(x)\mid p'(x)$，因此必有$p'(x)=0$.

\end{proof}

\begin{corollary}
若$\mathrm{ch}F=0$，则$F[x]$中任一不可约多项式在其分裂域中无重根。
\end{corollary}

\begin{proof}
设$p(x)$为$F[x]$中不可约多项式，则$\deg p'(x)=\deg p(x)-1\geqslant0$，故$p'(x)\neq0$。由上一推论，$p(x)$在其分裂域中只有单根.

\end{proof}

\begin{definition}
设$F$为域。
\begin{enumerate}[(1)]
\item 若$F[x]$中不可约多项式$p(x)$在其分裂域中只有单根，则称$p(x)$为$F$上\textbf{可分的不可约多项式}。
\item 若$F[x]$中多项式$f(x)$的每个不可约因式都是可分的，则称$f(x)$为$F$上\textbf{可分多项式}。
\end{enumerate}
\end{definition}

\begin{remark}
可分多项式具有如下性质：
\begin{enumerate}[(1)]
\item $F[x]$中不可约多项式$p(x)$可分当且仅当$p'(x)\neq0$.
\item 若$\mathrm{ch}F=0$，则$F[x]$的任何多项式都是可分的；仅当$\mathrm{ch}F=p>0$时，$F[x]$才可能存在不可分多项式。
\end{enumerate}
\end{remark}

\begin{theorem}
设$\mathrm{ch}F=p>0$，$f(x)$是$F[x]$中不可分不可约多项式，$K$是$f(x)$的分裂域，则在$K[x]$中有分解
\begin{align*}
f(x)=c(x-\alpha_1)^{p^e}(x-\alpha_2)^{p^e}\dots(x-\alpha_r)^{p^e},
\end{align*}
其中$i\neq j$时$\alpha_i\neq\alpha_j$，$e\in\mathbb{N}$，且
\begin{align*}
h(x)=c(x-\alpha_1^{p^e})(x-\alpha_2^{p^e})\dots(x-\alpha_r^{p^e})
\end{align*}
是$F[x]$中可分的不可约多项式，满足$f(x)=h(x^{p^e})$。
\end{theorem}

\begin{proof}
因$f(x)$不可分，故$f'(x)=0$。设$f(x)=\sum\limits_{k=0}^n a_kx^k$，则$ka_k=0,\forall 0\leqslant k\leqslant n$。
当$p\nmid k$时，必有$a_k=0$，于是
\begin{align*}
f(x)=a_{mp}x^{mp}+a_{(m-1)p}x^{(m-1)p}+\dots+a_px^p+a_0.
\end{align*}
令$g(x)=\sum\limits_{k=0}^m a_{kp}x^k$，则$g(x)\in F[x]$且$f(x)=g(x^p)$。
若$g(x)$可约，则$f(x)$可约，与$f(x)$不可约矛盾，故$g(x)$不可约。
重复上述过程，有限步后可得$F[x]$中不可约可分多项式$h(x)$，使得$f(x)=h(x^{p^e})$。

因$h(x)$可分，在其分裂域中有分解
\begin{align*}
h(x)=c(x-\beta_1)(x-\beta_2)\dots(x-\beta_r),
\end{align*}
其中$i\neq j$时$\beta_i\neq\beta_j$。于是
\begin{align*}
f(x)=c(x^{p^e}-\beta_1)(x^{p^e}-\beta_2)\dots(x^{p^e}-\beta_r).
\end{align*}
设$\alpha_i$是$x^{p^e}-\beta_i$的根，则$\beta_i=\alpha_i^{p^e}$，因此
\begin{align*}
f(x)=c(x-\alpha_1)^{p^e}(x-\alpha_2)^{p^e}\dots(x-\alpha_r)^{p^e}.
\end{align*}

\end{proof}

\begin{remark}
记$r=\deg h(x)$，即$f(x)$在其分裂域中不同根的个数，称为$f(x)$的\textbf{简约次数}，记作$\mathrm{red}f(x)$。
$\dfrac{\deg f(x)}{\mathrm{red}f(x)}$是素数$p$的方幂，且为$f(x)$的任一根的重数。
当且仅当$\deg f(x)=\mathrm{red}f(x)$时，$f(x)$为可分不可约多项式。
\end{remark}

\begin{definition}
设$F$为域，若$F[x]$中每个多项式都是可分多项式，则称$F$为\textbf{完备域}。
\end{definition}

\begin{theorem}
设$\mathrm{ch}F=p>0$，则$F$是完备域的充分必要条件是$\forall a\in F,\exists b\in F$使得$a=b^p$，即$F^p=\{a^p\mid a\in F\}=F$。
\end{theorem}

\begin{proof}
必要性：反证。若$F^p\neq F$，则存在$a\in F$满足$a\notin F^p$。令$f(x)=x^p-a$。
设$K$为$f(x)$的分裂域，$\alpha$为$f(x)$的根，则在$K[x]$中$f(x)=(x-\alpha)^p$。
若$f(x)$可约，则$f(x)=g(x)h(x)$，设$\deg g(x)=r$，则$g(x)=(x-\alpha)^r$，于是$\alpha^r\in F$。
因$(p,r)=1$，存在整数$u,v$使得$pu+vr=1$，故
\begin{align*}
\alpha=\alpha^{pu+vr}=a^u(\alpha^r)^v\in F,
\end{align*}
于是$a=\alpha^p\in F^p$，矛盾，故$f(x)$不可约。
又$f'(x)=0$，故$f(x)$不可分，因此$F$不是完备域。

充分性：若$F^p=F$，设$f(x)$为$F[x]$中不可分不可约多项式，则存在可分不可约多项式$h(x)\in F[x]$，使得$f(x)=h(x^{p^e}),e>1$。
设$h(x)=\sum\limits_{k=0}^m b_kx^k$，由$F^p=F$，存在$a_i\in F$使得$a_i^{p^e}=b_i$，于是
\begin{align*}
f(x)=h(x^{p^e})=\sum\limits_{k=0}^m a_k^{p^e}(x^k)^{p^e}=\left(\sum\limits_{k=0}^m a_kx^k\right)^{p^e},
\end{align*}
与$f(x)$不可约矛盾，故$F$是完备域.

\end{proof}

\begin{corollary}
有限域是完备域。
\end{corollary}

\begin{proof}
设$F$为有限域，$\mathrm{ch}F=p>0$，作映射$\sigma:F\to F,\sigma(a)=a^p$。
由$(ab)^p=a^pb^p$，$(a+b)^p=a^p+b^p$，知$\sigma$为域同态。
又$a^p=b^p\iff a=b$，故$\sigma$为单射。结合$F$有限，$\sigma$为满射，即$F^p=F$，因此$F$是完备域.

\end{proof}

\begin{corollary}
完备域$F$的代数扩张$K$也是完备域。
\end{corollary}

\begin{proof}
若$\mathrm{ch}F=0$，则$\mathrm{ch}K=0$，显然$K$是完备域。
设$\mathrm{ch}F=\mathrm{ch}K=p>0$，作映射$\sigma:K\to K,\sigma(a)=a^p$。
任取$\alpha\in K$，令$E=F(\alpha)$。由$F$完备，$\sigma(F)=F$。
$\sigma(E)\cong E$，且$[\sigma(E):\sigma(F)]=[E:F]$，故$\sigma(E)=E$，即$E^p=E$。
因此存在$\beta\in E\subseteq K$使得$\alpha=\beta^p$，即$K^p=K$，故$K$是完备域.

\end{proof}

\begin{corollary}
设$\mathbb{Z}_p(t)$是$\mathbb{Z}_p$的单超越扩张，则$\mathbb{Z}_p(t)$不是完备域。
\end{corollary}

\begin{proof}
$x^p-t\in\mathbb{Z}_p(t)[x]$是不可约多项式，且$(x^p-t)'=0$，故$x^p-t$是不可分不可约多项式，因此$\mathbb{Z}_p(t)$不是完备域.

\end{proof}





\end{document}